% Created 2026-09-14 lun 21:08
% Intended LaTeX compiler: pdflatex
\documentclass[11pt]{article}

\usepackage[utf8]{inputenc}
\usepackage[T1]{fontenc}
\usepackage{graphicx}
\usepackage{longtable}
\usepackage{wrapfig}
\usepackage{rotating}
\usepackage[normalem]{ulem}
\usepackage{amsmath}
\usepackage{amssymb}
\usepackage{capt-of}
\usepackage{hyperref}
\usepackage{color}
\usepackage{listings}
\author{alvaro}
\date{\today}
\title{}
\hypersetup{
 pdfauthor={alvaro},
 pdftitle={},
 pdfkeywords={},
 pdfsubject={},
 pdfcreator={},
 pdflang={English}}
\begin{document}

\tableofcontents

\section{Descripción general}
\label{sec:org000000f}

\subsection{Arquitectura final}
\label{sec:org0000000}

\begin{center}
\includesvg[width=.9\linewidth]{./media/arquitectura}
\label{}
\end{center}
\subsection{Redes Docker}
\label{sec:org0000003}
\begin{itemize}
\item \texttt{red-publica} (bridge) con subred \texttt{172.21.0.0/24}: contiene los servicios expuestos al exterior: \texttt{dns}, \texttt{proxy} y \texttt{vpn}.
\item \texttt{red-interna} (bridge) con subred \texttt{172.22.0.0/24}: contiene los servicios \emph{backend}: \texttt{apache}, \texttt{db} y \texttt{phpmyadmin}. \texttt{proxy} y \texttt{vpn} también se conectan a esta red para alcanzar el backend.
\item Dentro de una misma red Docker los contenedores se resuelven entre sí por su \textbf{nombre de servicio} (\texttt{db}, \texttt{apache}\ldots{}). Los nombres públicos del dominio \texttt{asir.test} los resuelve el servidor DNS del proyecto.
\end{itemize}
\subsection{Contenedores e IPs}
\label{sec:org0000006}

\begin{center}
\begin{tabular}{lllrl}
Contenedor & Imagen & IP red-publica & IP red-interna & Funciones principales\\
\hline
\texttt{dns} & bind9 & 172.21.0.53 & — & Autoritativo de \texttt{asir.test} y recursivo para el resto\\
\texttt{proxy} & nginx & 172.21.0.10 & 172.22.0.10 & Proxy inverso a Apache y a PHPMyAdmin\\
\texttt{apache} & php:8.2-apache & — & 172.22.0.11 & Apache + PHP con WordPress, SSH (fase 5)\\
\texttt{db} & mysql:8.4 & — & 172.22.0.12 & MySQL con la base de datos de WordPress, SSH (fase 5)\\
\texttt{phpmyadmin} & phpmyadmin & — & 172.22.0.13 & Gestión web de \texttt{db}, accesible solo por proxy o por VPN\\
\texttt{vpn} & wireguard (wg-easy) & 172.21.0.20 & 172.22.0.20 & Servidor WireGuard con panel web de gestión\\
\end{tabular}
\end{center}
\subsection{Puertos expuestos en el host}
\label{sec:org0000009}

\begin{center}
\begin{tabular}{lll}
Puerto & Servicio & Uso\\
\hline
53 (TCP+UDP) & \texttt{dns} & Consultas DNS de clientes externos y de clientes VPN\\
80 y 443 & \texttt{proxy} & WordPress (\texttt{www.asir.test}) y PHPMyAdmin (\texttt{pma.asir.test})\\
51820/udp & \texttt{vpn} & Túnel WireGuard para los clientes VPN\\
51821/tcp & \texttt{vpn} & Panel web de gestión de clientes VPN (solo administrador)\\
\end{tabular}
\end{center}
\subsection{Dominio}
\label{sec:org000000c}
Se emplea el dominio reservado para pruebas \texttt{asir.test} (RFC 6761). Registros requeridos:
\begin{itemize}
\item \texttt{dns.asir.test}
\item \texttt{proxy.asir.test}
\item \texttt{vpn.asir.test}
\item \texttt{apache.asir.test}
\item \texttt{db.asir.test}
\item \texttt{pma.asir.test},
\item alias \texttt{phpmyadmin.asir.test} (CNAME de \texttt{pma})
\item Alias \texttt{www.asir.test} (CNAME de \texttt{proxy}).
\end{itemize}
\section{Fase 1 — Servidor LAMP}
\label{sec:org0000021}

\subsection{Enunciado}
\label{sec:org0000012}
\begin{itemize}
\item Desplegar un servidor LAMP (Linux + Apache + MySQL + PHP) en \emph{tres contenedores}:
\texttt{apache} (Apache con PHP), \texttt{db} (MySQL) y \texttt{phpmyadmin} (gestión de \texttt{db}).
\item Los tres contenedores deben formar parte de una red Docker propia (\texttt{red-interna})
y comunicarse entre sí por sus nombres de servicio.
\item Debe probarse una página PHP que muestre el resultado de una consulta a MySQL.
\end{itemize}
\subsection{Objetivos}
\label{sec:org0000015}
\begin{itemize}
\item Crear y administrar redes Docker (driver \texttt{bridge}, subred propia, IPs fijas).
\item Conocer las variables de entorno de las imágenes MySQL y PHPMyAdmin.
\item Conectar un contenedor PHP a la base de datos usando únicamente el nombre \texttt{db}.
\end{itemize}
\subsection{Tareas}
\label{sec:org0000018}
\begin{enumerate}
\item Crear el fichero \texttt{docker-compose.yml} con los servicios \texttt{db}, \texttt{apache} y
\texttt{phpmyadmin} y la red \texttt{red-interna} (172.22.0.0/24) con IPs fijas.
\item Configurar \texttt{db} con: volumen persistente (\texttt{./db/data}), base de datos
\texttt{wordpress}, usuario \texttt{wp\_user} y contraseñas.
\item Configurar \texttt{apache} con la imagen \texttt{php:8.2-apache} y montar el volumen
\texttt{./apache/www} en \texttt{/var/www/html}.
\item Configurar \texttt{phpmyadmin} apuntando a \texttt{PMA\_HOST=db}.
\item Añadir temporalmente \texttt{8080:80} (apache) y \texttt{8081:80} (phpmyadmin) para poder
probarlos desde el host. Se retirarán en la fase 4.
\item Crear en \texttt{apache/www} los ficheros \texttt{index.php} (bienvenida del LAMP) y
\texttt{test-db.php} (consulta a MySQL con \texttt{mysqli}).
\end{enumerate}
\subsection{Comprobaciones}
\label{sec:org000001b}
\begin{itemize}
\item \texttt{docker compose ps} muestra los tres contenedores arriba.
\item Acceder a \texttt{http://localhost:8080/index.php} y \texttt{http://localhost:8080/test-db.php}.
\item \texttt{http://localhost:8081} muestra el login de PHPMyAdmin; se entra con
\texttt{root} / \texttt{asir\_root} y la base \texttt{wordpress} debe existir.
\item Desde \texttt{apache} debe resolverse el nombre \texttt{db}:
\texttt{docker compose exec apache getent hosts db}.
\end{itemize}

\begin{lstlisting}[language=bash,numbers=none]
curl http://localhost:8080/test-db.php
curl -u root:asir_root http://localhost:8081/
\end{lstlisting}
\subsection{Entrega}
\label{sec:org000001e}
\begin{itemize}
\item \texttt{docker-compose.yml}, \texttt{index.php} y \texttt{test-db.php}.
\item Capturas de la resolución de \texttt{db} y de las dos páginas web funcionando.
\end{itemize}
\section{Fase 2 — Servidor DNS}
\label{sec:org0000033}

\subsection{Enunciado}
\label{sec:org0000024}
\begin{itemize}
\item Desplegar un contenedor \texttt{dns} con BIND9 que sea \textbf{autoritativo} para el dominio
\texttt{asir.test} y \textbf{recursivo} para el resto de nombres.
\item Debe almacenar un registro de nombre por cada contenedor del proyecto y los
registros de resolución inversa (PTR).
\item Los clientes externos deben poder utilizarlo: desde el host se apunta el
\emph{resolver} a \texttt{172.21.0.53}.
\end{itemize}
\subsection{Objetivos}
\label{sec:org0000027}
\begin{itemize}
\item Montar BIND9 en un contenedor con ficheros de configuración propios.
\item Crear zona directa e inversas con registros A, CNAME y PTR.
\item Configurar la recursividad y comprobar con \texttt{dig} / \texttt{nslookup}.
\end{itemize}
\subsection{Tareas}
\label{sec:org000002a}
\begin{enumerate}
\item Añadir el servicio \texttt{dns} con la imagen \texttt{internetsystemsconsortium/bind9} y
publicar los puertos \texttt{53} (TCP y UDP).
\item Crear la red \texttt{red-publica} (172.21.0.0/24) y asignar a \texttt{dns} la IP fija
172.21.0.53.
\item Montar en el contenedor \texttt{named.conf}, \texttt{named.conf.options}, \texttt{named.conf.local}
y el directorio \texttt{zones/}.
\item Definir la zona directa \texttt{asir.test} con un A por contenedor (\texttt{dns}, \texttt{proxy},
\texttt{vpn}, \texttt{apache}, \texttt{db}, \texttt{pma}) y los CNAME \texttt{phpmyadmin} → \texttt{pma} y \texttt{www} → \texttt{proxy}.
\item Definir las zonas inversas \texttt{21.172.in-addr.arpa} y \texttt{22.172.in-addr.arpa} con
los PTR de todas las IPs del proyecto.
\item Configurar la recursividad con \texttt{forwarders} (p. ej. 1.1.1.1 y 8.8.8.8).
\end{enumerate}
\subsection{Comprobaciones}
\label{sec:org000002d}
\begin{itemize}
\item Desde el host, apuntando el \emph{resolver} a 172.21.0.53:
\end{itemize}

\begin{lstlisting}[language=bash,numbers=none]
dig @172.21.0.53 www.asir.test A
dig @172.21.0.53 apache.asir.test A
dig @172.21.0.53 -x 172.22.0.12
dig @172.21.0.53 -x 172.21.0.53
dig @172.21.0.53 google.es A    # recursividad hacia el exterior
\end{lstlisting}
\begin{itemize}
\item Todos los nombres devuelven el registro esperado; las inversas devuelven el
nombre correspondiente.
\end{itemize}
\subsection{Entrega}
\label{sec:org0000030}
\begin{itemize}
\item Configuración de BIND9 y ficheros de zona.
\item Salida de \texttt{dig} de al menos 4 nombres directos, 2 inversos y 1 recursivo.
\end{itemize}
\section{Fase 3 — WordPress instalado en LAMP}
\label{sec:org0000045}

\subsection{Enunciado}
\label{sec:org0000036}
\begin{itemize}
\item Instalar WordPress sobre el LAMP de la fase 1, de modo que use la base de datos
\texttt{db} y el servidor web \texttt{apache}.
\item Completar la instalación por web y crear al menos una entrada (página o artículo).
\end{itemize}
\subsection{Objetivos}
\label{sec:org0000039}
\begin{itemize}
\item Automatizar la descarga de WordPress mediante un \emph{script} (\texttt{scripts/install.sh}).
\item Generar \texttt{wp-config.php} con las credenciales de \texttt{db} y las \textbf{salts}.
\item Comprobar que una aplicación PHP real persiste sus datos en MySQL.
\end{itemize}
\subsection{Tareas}
\label{sec:org000003c}
\begin{enumerate}
\item En \texttt{scripts/install.sh}: descargar \texttt{https://wordpress.org/latest.tar.gz} y
extraerlo en \texttt{apache/www}. No sobrescribir si ya existe \texttt{wp-config.php}.
\item Generar \texttt{wp-config.php} con \texttt{DB\_NAME=wordpress}, \texttt{DB\_USER=wp\_user},
\texttt{DB\_PASSWORD=wp\_pass\_2026}, \texttt{DB\_HOST=db} y las 8 \emph{salts} generadas con
\texttt{openssl rand -base64 48}.
\item Levantar el escenario y completar el asistente web apuntando el navegador a
\texttt{http://localhost:8080} (en esta fase todavía sin proxy).
\item Crear un artículo de prueba y comprobar que tras
\texttt{docker compose restart db} el artículo sigue existiendo (persistencia en
\texttt{db/data}).
\end{enumerate}
\subsection{Comprobaciones}
\label{sec:org000003f}
\begin{itemize}
\item \#+begin\textsubscript{src} bash
docker compose exec db mysql -u wp\textsubscript{user} -pwp\textsubscript{pass}\textsubscript{2026} wordpress -e "show tables;"
\#+end\textsubscript{src}
muestra las tablas de WordPress.
\item \texttt{/wp-admin/} permite iniciar sesión con el usuario creado en el asistente.
\end{itemize}
\subsection{Entrega}
\label{sec:org0000042}
\begin{itemize}
\item Fragmento relevante de \texttt{install.sh}, \texttt{wp-config.php} y captura del sitio ya
instalado.
\end{itemize}
\section{Fase 4 — Proxy inverso}
\label{sec:org0000057}

\subsection{Enunciado}
\label{sec:org0000048}
\begin{itemize}
\item Colocar un contenedor \texttt{proxy} (nginx) como puerta de entrada única al servidor
Apache y al servidor PHPMyAdmin.
\item El proxy usará dos nombres virtuales: \texttt{www.asir.test} → \texttt{apache:80} y
\texttt{pma.asir.test} → \texttt{phpmyadmin:80}.
\item A partir de esta fase los clientes acceden únicamente a través del proxy; se
retiran las publicaciones \texttt{8080} y \texttt{8081}.
\end{itemize}
\subsection{Objetivos}
\label{sec:org000004b}
\begin{itemize}
\item Comprender el patrón de proxy inverso y la directiva \texttt{proxy\_pass}.
\item Aplicar enrutamiento por \texttt{server\_name}.
\item Ocultar los servidores backend de la red exterior.
\end{itemize}
\subsection{Tareas}
\label{sec:org000004e}
\begin{enumerate}
\item Añadir el servicio \texttt{proxy} (imagen \texttt{nginx:1.27-alpine}) a \texttt{red-publica}
(172.21.0.10) y a \texttt{red-interna} (172.22.0.10).
\item Montar \texttt{./proxy/conf.d/proxy.conf} en \texttt{/etc/nginx/conf.d}.
\item Publicar solo los puertos \texttt{80} y \texttt{443} de \texttt{proxy} en el host.
\item Retirar las publicaciones \texttt{8080} y \texttt{8081} de \texttt{apache} y \texttt{phpmyadmin}.
\item Definir los dos bloques \texttt{server\{\}:} para \texttt{www.asir.test} y \texttt{pma.asir.test} con
\texttt{proxy\_pass http://apache:80} y \texttt{proxy\_pass http://phpmyadmin:80}
respectivamente, propagando \texttt{Host} y \texttt{X-Forwarded-For}.
\item Comprobar desde el host añadiendo una entrada estática al DNS del sistema
(o usando \texttt{-{}-{}resolve} con \texttt{curl}) y después mediante el DNS de la fase 2.
\end{enumerate}
\subsection{Comprobaciones}
\label{sec:org0000051}
\begin{itemize}
\item \#+begin\textsubscript{src} bash
curl -s -o \emph{dev/null -w "\%\{http\textsubscript{code}\}\n" -H "Host: www.asir.test" \url{http://127.0.0.1}}
curl -s -o \emph{dev/null -w "\%\{http\textsubscript{code}\}\n" -H "Host: pma.asir.test" \url{http://127.0.0.1}}
\#+end\textsubscript{src}
\item Abriendo \texttt{http://www.asir.test} y \texttt{http://pma.asir.test} en el navegador (con
DNS del sistema apuntando a 172.21.0.53) se ven WordPress y PHPMyAdmin.
\end{itemize}
\subsection{Entrega}
\label{sec:org0000054}
\begin{itemize}
\item \texttt{proxy.conf} completo y capturas de ambos sitios funcionando a través del proxy.
\end{itemize}
\section{Fase 5 — Acceso por VPN WireGuard}
\label{sec:org0000069}

\subsection{Enunciado}
\label{sec:org000005a}
\begin{itemize}
\item Desplegar un servidor VPN WireGuard (\texttt{vpn}).
\item Los clientes conectados a la VPN deben:
\begin{itemize}
\item usar el servidor DNS del proyecto (172.21.0.53) para resolver \texttt{asir.test};
\item acceder al \emph{backend} PHPMyAdmin (\texttt{pma.asir.test}) directamente, sin pasar por
el proxy;
\item conectarse por SSH a los contenedores \texttt{apache} y \texttt{db}.
\end{itemize}
\end{itemize}
\subsection{Objetivos}
\label{sec:org000005d}
\begin{itemize}
\item Instalar y administrar WireGuard con wg-easy (panel web para clientes).
\item Encaminar el tráfico de la VPN hacia las subredes de Docker.
\item Publicar servicios del backend de forma segura a través del túnel.
\end{itemize}
\subsection{Tareas}
\label{sec:org0000060}
\begin{enumerate}
\item Añadir el servicio \texttt{vpn} con la imagen \texttt{weejewel/wg-easy}, \texttt{cap\_add NET\_ADMIN},
\texttt{net.ipv4.ip\_forward} y \texttt{net.ipv4.conf.all.src\_valid\_mark=1}.
\item Publicar \texttt{51820/udp} (túnel) y \texttt{51821/tcp} (panel web de gestión).
\item Conectar \texttt{vpn} a \texttt{red-publica} (172.21.0.20) y a \texttt{red-interna} (172.22.0.20)
para alcanzar el backend.
\item Configurar las variables del entorno: \texttt{WG\_HOST} (IP del host o del servidor),
\texttt{PASSWORD}, \texttt{WG\_DEFAULT\_DNS=172.21.0.53} y \texttt{WG\_ALLOWED\_IPS} con las subredes
públicas, interna y la del túnel (172.21.0.0/24, 172.22.0.0/24 y 10.9.0.0/24).
\item Instalar \texttt{openssh-server} en \texttt{apache} y \texttt{db} (permitiendo login de root con
contraseña) e iniciarlo junto con el servicio principal.
\item Crear un cliente VPN desde el panel \texttt{http://IP\_HOST:51821} y descargar su
configuración.
\end{enumerate}
\subsection{Comprobaciones (desde el cliente VPN conectado)}
\label{sec:org0000063}
\begin{itemize}
\item \#+begin\textsubscript{src} bash
ip addr show wg0                     \# IP 10.9.0.x asignada
ping -c 3 172.22.0.12                \# alcanza MySQL por la VPN
dig apache.asir.test @172.21.0.53    \# usa el DNS del proyecto
curl \url{http://pma.asir.test/}           \# backend PHPMyAdmin por la VPN
ssh root@apache.asir.test            \# shell en el contenedor apache
ssh root@db.asir.test                \# shell en el contenedor db
\#+end\textsubscript{src}
\item En el host se comprueba que el puerto 8081 (PHPMyAdmin) NO está publicado:
solo se alcanza dentro de la VPN o a través de \texttt{pma.asir.test} (proxy).
\end{itemize}
\subsection{Entrega}
\label{sec:org0000066}
\begin{itemize}
\item Configuración del servicio \texttt{vpn} y capturas del cliente VPN con todos los
apartados anteriores funcionando.
\end{itemize}
\section{Evaluación}
\label{sec:org000006c}

\begin{center}
\begin{tabular}{ll}
Criterio & Peso\\
\hline
Fase 1: LAMP operativo y explicación de sus componentes & 15\%\\
Fase 2: resolución directa e inversa de todos los contenedores & 20\%\\
Fase 3: WordPress instalado y con persistencia en MySQL & 15\%\\
Fase 4: proxy inverso con doble nombre virtual y backend oculto & 20\%\\
Fase 5: cliente VPN con DNS + accesso al backend y SSH a apache/db & 20\%\\
Despliegue automático con un solo \texttt{docker compose} y los scripts & 10\%\\
\end{tabular}
\end{center}
\section{Solución}
\label{sec:org00000a5}

\subsection{Estructura definitiva del repositorio}
\label{sec:org000006f}

\begin{lstlisting}[language=text,numbers=none]
proyecto-asir/
├── docker-compose.yml
├── .env.example
├── apache/
│   ├── Dockerfile
│   ├── entrypoint.sh
│   └── www/                 # WordPress (lo descarga scripts/install.sh)
├── db/
│   ├── Dockerfile
│   ├── entrypoint-ssh.sh
│   └── data/
├── proxy/
│   └── conf.d/proxy.conf
├── dns/
│   ├── named.conf
│   ├── named.conf.options
│   ├── named.conf.local
│   └── zones/
│       ├── db.asir.test
│       ├── db.172.21
│       └── db.172.22
└── scripts/
    ├── install.sh
    └── comprobaciones.sh
\end{lstlisting}
\subsection{docker-compose.yml}
\label{sec:org0000072}

\begin{lstlisting}[language=yaml,numbers=none]
name: proyecto-asir

services:

  proxy:
    image: nginx:1.27-alpine
    container_name: proxy
    restart: unless-stopped
    ports:
      - "80:80"
      - "443:443"
    volumes:
      - ./proxy/conf.d:/etc/nginx/conf.d:ro
    depends_on:
      - apache
      - phpmyadmin
    networks:
      red-publica:
        ipv4_address: 172.21.0.10
      red-interna:
        ipv4_address: 172.22.0.10

  apache:
    build: ./apache
    container_name: apache
    restart: unless-stopped
    volumes:
      - ./apache/www:/var/www/html
    depends_on:
      - db
    networks:
      red-interna:
        ipv4_address: 172.22.0.11

  db:
    build: ./db
    container_name: db
    restart: unless-stopped
    environment:
      MYSQL_ROOT_PASSWORD: asir_root
      MYSQL_DATABASE: wordpress
      MYSQL_USER: wp_user
      MYSQL_PASSWORD: wp_pass_2026
    volumes:
      - ./db/data:/var/lib/mysql
    networks:
      red-interna:
        ipv4_address: 172.22.0.12

  phpmyadmin:
    image: phpmyadmin:5.2
    container_name: phpmyadmin
    restart: unless-stopped
    environment:
      PMA_HOST: db
      PMA_PORT: 3306
      PMA_USER: root
      PMA_PASSWORD: asir_root
      PMA_ABSOLUTE_URI: "http://pma.asir.test"
    depends_on:
      - db
    networks:
      red-interna:
        ipv4_address: 172.22.0.13

  dns:
    image: internetsystemsconsortium/bind9:9.19
    container_name: dns
    restart: unless-stopped
    ports:
      - "53:53/udp"
      - "53:53/tcp"
    volumes:
      - ./dns/named.conf:/etc/bind/named.conf:ro
      - ./dns/named.conf.options:/etc/bind/named.conf.options:ro
      - ./dns/named.conf.local:/etc/bind/named.conf.local:ro
      - ./dns/zones:/etc/bind/zones:ro
    networks:
      red-publica:
        ipv4_address: 172.21.0.53

  vpn:
    image: weejewel/wg-easy
    container_name: vpn
    restart: unless-stopped
    cap_add:
      - NET_ADMIN
      - SYS_MODULE
    sysctls:
      - net.ipv4.conf.all.src_valid_mark=1
      - net.ipv4.ip_forward=1
    ports:
      - "51820:51820/udp"
      - "51821:51821/tcp"
    environment:
      WG_HOST: "${WG_HOST:-127.0.0.1}"
      PASSWORD: "${WG_PASSWORD:-asir-vpn-2026}"
      WG_DEFAULT_ADDRESS: "10.9.0.x"
      WG_DEFAULT_DNS: "172.21.0.53"
      WG_ALLOWED_IPS: "172.21.0.0/24, 172.22.0.0/24, 10.9.0.0/24"
      WG_PERSISTENT_KEEPALIVE: "25"
    networks:
      red-publica:
        ipv4_address: 172.21.0.20
      red-interna:
        ipv4_address: 172.22.0.20

networks:
  red-publica:
    driver: bridge
    ipam:
      config:
        - subnet: 172.21.0.0/24
  red-interna:
    driver: bridge
    ipam:
      config:
        - subnet: 172.22.0.0/24
\end{lstlisting}
\subsection{.env.example}
\label{sec:org0000075}

\begin{lstlisting}[language=text,numbers=none]
# IP/dominio que usan los clientes VPN para conectarse (IP del host o publica)
WG_HOST=192.168.1.10
# Contrasena del panel web de WireGuard (http://IP_HOST:51821)
WG_PASSWORD=asir-vpn-2026
\end{lstlisting}
\subsection{apache/Dockerfile}
\label{sec:org0000078}

\begin{lstlisting}[language=dockerfile,numbers=none]
FROM php:8.2-apache

RUN apt-get update \
    && apt-get install -y --no-install-recommends vim openssh-server \
    && docker-php-ext-install mysqli pdo_mysql \
    && a2enmod rewrite \
    && sed -i 's/^#PermitRootLogin.*/PermitRootLogin yes/' /etc/ssh/sshd_config \
    && mkdir -p /run/sshd \
    && echo 'root:asir' | chpasswd

COPY entrypoint.sh /usr/local/bin/entrypoint.sh
RUN chmod +x /usr/local/bin/entrypoint.sh

EXPOSE 80 22
ENTRYPOINT ["/usr/local/bin/entrypoint.sh"]
\end{lstlisting}
\subsection{apache/entrypoint.sh}
\label{sec:org000007b}

\begin{lstlisting}[language=bash,numbers=none]
#!/bin/bash
/usr/sbin/sshd
exec apache2-foreground
\end{lstlisting}
\subsection{db/Dockerfile}
\label{sec:org000007e}

\begin{lstlisting}[language=dockerfile,numbers=none]
FROM mysql:8.4

RUN apt-get update \
    && apt-get install -y --no-install-recommends openssh-server vim \
    && sed -i 's/^#PermitRootLogin.*/PermitRootLogin yes/' /etc/ssh/sshd_config \
    && mkdir -p /run/sshd \
    && echo 'root:asir' | chpasswd

COPY entrypoint-ssh.sh /usr/local/bin/entrypoint-ssh.sh
RUN chmod +x /usr/local/bin/entrypoint-ssh.sh

EXPOSE 3306 22
ENTRYPOINT ["/usr/local/bin/entrypoint-ssh.sh"]
CMD ["mysqld"]
\end{lstlisting}
\subsection{db/entrypoint-ssh.sh}
\label{sec:org0000081}

\begin{lstlisting}[language=bash,numbers=none]
#!/bin/bash
/usr/sbin/sshd
exec docker-entrypoint.sh "$@"
\end{lstlisting}
\subsection{proxy/conf.d/proxy.conf}
\label{sec:org0000084}

\begin{lstlisting}[language=nginx,numbers=none]
# WordPress a traves del proxy
server {
    listen 80;
    server_name www.asir.test;

    location / {
        proxy_pass http://apache:80;
        proxy_set_header Host $host;
        proxy_set_header X-Real-IP $remote_addr;
        proxy_set_header X-Forwarded-For $proxy_add_x_forwarded_for;
    }
}

# PHPMyAdmin a traves del proxy
server {
    listen 80;
    server_name pma.asir.test phpmyadmin.asir.test;

    location / {
        proxy_pass http://phpmyadmin:80;
        proxy_set_header Host $host;
        proxy_set_header X-Real-IP $remote_addr;
        proxy_set_header X-Forwarded-For $proxy_add_x_forwarded_for;
    }
}
\end{lstlisting}
\subsection{dns/named.conf}
\label{sec:org0000087}

\begin{lstlisting}[language=conf,numbers=none]
include "/etc/bind/named.conf.options";
include "/etc/bind/named.conf.local";
\end{lstlisting}
\subsection{dns/named.conf.options}
\label{sec:org000008a}

\begin{lstlisting}[language=conf,numbers=none]
options {
    directory "/var/cache/bind";

    recursion yes;
    allow-query { any; };

    listen-on port 53  { any; };
    listen-on-v6       { any; };

    forwarders { 1.1.1.1; 8.8.8.8; };
    dnssec-validation no;
};
\end{lstlisting}
\subsection{dns/named.conf.local}
\label{sec:org000008d}

\begin{lstlisting}[language=conf,numbers=none]
// Zona directa del proyecto
zone "asir.test" {
    type master;
    file "/etc/bind/zones/db.asir.test";
};

// Zonas inversas de las dos redes Docker
zone "21.172.in-addr.arpa" {
    type master;
    file "/etc/bind/zones/db.172.21";
};

zone "22.172.in-addr.arpa" {
    type master;
    file "/etc/bind/zones/db.172.22";
};
\end{lstlisting}
\subsection{dns/zones/db.asir.test}
\label{sec:org0000090}

\begin{lstlisting}[language=conf,numbers=none]
$TTL 86400
@   IN SOA dns.asir.test. admin.asir.test. (
        2026091401 ; serial
        3600       ; refresh
        1800       ; retry
        604800     ; expire
        86400 )    ; minimum

    IN NS  dns.asir.test.

; Servidores
dns     IN A    172.21.0.53
proxy   IN A    172.21.0.10
vpn     IN A    172.21.0.20

; Backend (solo red-interna / VPN)
apache  IN A    172.22.0.11
db      IN A    172.22.0.12
pma     IN A    172.22.0.13

; Alias
phpmyadmin  IN CNAME pma
www         IN CNAME proxy
\end{lstlisting}
\subsection{dns/zones/db.172.21}
\label{sec:org0000093}

\begin{lstlisting}[language=conf,numbers=none]
$TTL 86400
@   IN SOA dns.asir.test. admin.asir.test. (
        2026091401 ; serial
        3600       ; refresh
        1800       ; retry
        604800     ; expire
        86400 )    ; minimum

    IN NS  dns.asir.test.

10  IN PTR proxy.asir.test.
20  IN PTR vpn.asir.test.
53  IN PTR dns.asir.test.
\end{lstlisting}
\subsection{dns/zones/db.172.22}
\label{sec:org0000096}

\begin{lstlisting}[language=conf,numbers=none]
$TTL 86400
@   IN SOA dns.asir.test. admin.asir.test. (
        2026091401 ; serial
        3600       ; refresh
        1800       ; retry
        604800     ; expire
        86400 )    ; minimum

    IN NS  dns.asir.test.

10  IN PTR proxy.asir.test.
11  IN PTR apache.asir.test.
12  IN PTR db.asir.test.
13  IN PTR pma.asir.test.
\end{lstlisting}
\subsection{scripts/install.sh}
\label{sec:org0000099}

\begin{lstlisting}[language=bash,numbers=none]
#!/bin/bash
# Preparacion automatica del proyecto ASIR.
# Uso: ./scripts/install.sh
set -euo pipefail
cd "$(dirname "$0")/.."

echo "[1/3] Creando directorios de datos"
mkdir -p apache/www db/data

echo "[2/3] Desplegando WordPress (fases 1 y 3)"
if [ ! -f "apache/www/wp-config.php" ]; then
  curl -fsSL https://wordpress.org/latest.tar.gz -o /tmp/wordpress.tar.gz
  tar -xzf /tmp/wordpress.tar.gz -C apache/www --strip-components=1

  SALTS=""
  for i in AUTH_KEY SECURE_AUTH_KEY LOGGED_IN_KEY NONCE_KEY \
           AUTH_SALT SECURE_AUTH_SALT LOGGED_IN_SALT NONCE_SALT; do
    SALTS+=$(printf "define('%s', '%s');\n" "$i" "$(openssl rand -base64 48)")
  done

  cat > apache/www/wp-config.php <<EOF
<?php
define('DB_NAME', 'wordpress');
define('DB_USER', 'wp_user');
define('DB_PASSWORD', 'wp_pass_2026');
define('DB_HOST', 'db');
define('DB_CHARSET', 'utf8mb4');
define('DB_COLLATE', '');

$SALTS
\$table_prefix = 'wp_';

define('WP_DEBUG', false);

// Acceso final a traves del proxy inverso
define('WP_HOME', 'http://www.asir.test');
define('WP_SITEURL', 'http://www.asir.test');

if ( ! defined( 'ABSPATH' ) ) {
    define( 'ABSPATH', __DIR__ . '/' );
}
require_once ABSPATH . 'wp-settings.php';
EOF
  echo "  - wp-config.php generado"
fi

echo "[3/3] Configurando WireGuard (fase 5)"
if [ ! -f .env ]; then
  HOST_IP=$(hostname -I | awk '{print $1}')
  cat > .env <<EOF
WG_HOST=${HOST_IP}
WG_PASSWORD=asir-vpn-2026
EOF
  echo "  - .env generado (WG_HOST=${HOST_IP})"
fi

echo "Listo. Ejecuta:  docker compose up -d --build"
\end{lstlisting}
\subsection{scripts/comprobaciones.sh}
\label{sec:org000009c}

\begin{lstlisting}[language=bash,numbers=none]
#!/bin/bash
# Verificacion del despliegue completo.  set -euo pipefail cd
"$(dirname "$0")/.."

DNS_IP=172.21.0.53

echo "=== 1. Contenedores ===" docker compose ps --format "table
{{.Name}}\t{{.Status}}\t{{.Ports}}"

echo "=== 2. Resolucion directa asir.test ===" for n in dns proxy www
vpn apache db pma phpmyadmin; do r=$(dig +short "@${DNS_IP}"
"${n}.asir.test" A | head -n1) echo " ${n}.asir.test -> ${r:-SIN
RESPUESTA}" done

echo "=== 3. Resolucion inversa ===" for ip in 172.21.0.53 172.21.0.10
172.21.0.20 \ 172.22.0.11 172.22.0.12 172.22.0.13; do echo " ${ip} ->
$(dig +short -x "${ip}" "@${DNS_IP}")" done

echo "=== 4. Proxy inverso ===" curl -s -o /dev/null -w "
www.asir.test -> HTTP %{http_code}\n" \ -H "Host: www.asir.test"
http://127.0.0.1/ curl -s -o /dev/null -w " pma.asir.test -> HTTP
%{http_code}\n" \ -H "Host: pma.asir.test" http://127.0.0.1/

echo "=== 5. SSH en backend (fase 5) ===" docker compose exec apache
bash -c 'pgrep sshd >/dev/null && echo " sshd OK en apache"' docker
compose exec db bash -c 'pgrep sshd >/dev/null && echo " sshd OK en
 db"'

echo "=== 6. Log de WireGuard ===" docker compose logs vpn 2>/dev/null
| tail -n 3
\end{lstlisting}
\subsection{Puesta en marcha}
\label{sec:org000009f}

\begin{lstlisting}[language=bash,numbers=none]
chmod +x scripts/*.sh ./scripts/install.sh # descarga WordPress,
genera wp-config y .env docker compose up -d --build
\end{lstlisting}

\begin{itemize}
\item Acceder a \texttt{http://www.asir.test} (WordPress) y \texttt{http://pma.asir.test} (PHPMyAdmin) con el DNS del sistema apuntando a 172.21.0.53.
\item Crear un cliente VPN en \texttt{http://IP\_DEL\_HOST:51821} y descargar el fichero \texttt{*.conf}. Importarlo en el cliente WireGuard, conectarse y repetir la lista de comprobaciones de la fase 5.
\end{itemize}
\subsection{Notas sobre wg-easy}
\label{sec:org00000a2}
\begin{itemize}
\item \texttt{WG\_HOST} debe ser la IP o dominio que el cliente VPN usa para conectarse: en prácticas de aula basta la IP del host (se detecta en \texttt{install.sh}); con clientes externos hay que usar la IP pública o un dominio con registro A.
\item Si se prefiere autenticar con hash en lugar de \texttt{PASSWORD}, se genera con \texttt{docker run -{}-{}rm weejewel/wg-easy wgpw 'asir-vpn-2026'} y se asigna la variable \texttt{PASSWORD\_HASH} en el servicio \texttt{vpn}.
\item wg-easy asigna IPs automáticas dentro de \texttt{10.9.0.0/24} a cada cliente creado desde el panel.
\end{itemize}
\end{document}
